% Related Work

\label{sec:related-work}

This section positions the project relative to the baselines used in the interim experiments and the GRPO motivation for the final project stage. The aim is not to provide a broad survey of reinforcement learning, but to justify the comparison against PPO, SAC, and TD3 and to identify the gap addressed by a later GRPO-control variant. These baselines represent three central families in deep reinforcement learning for control: on-policy policy optimization, off-policy stochastic actor-critic learning, and deterministic actor-critic learning for continuous action spaces.

PPO is the natural on-policy baseline for this project. It optimizes a policy using sampled trajectories and stabilizes policy-gradient updates through a clipped surrogate objective, limiting how far the updated policy can move from the previous policy in a single update \citep{schulman2017proximal}. PPO is also closely connected to advantage estimation. Generalized Advantage Estimation provides a practical bias--variance trade-off for estimating the advantage signals used in policy updates \citep{schulman2015gae}. In this report, PPO therefore serves both as an established control baseline and as the closest reference point for the later GRPO discussion.

SAC and TD3 complement PPO by representing off-policy actor-critic approaches. SAC uses a maximum-entropy objective, optimizing both expected return and policy entropy \citep{haarnoja2018sacapps}. This makes SAC a relevant stochastic off-policy baseline, especially in control tasks where exploration and reuse of past experience are important. TD3 instead follows the deterministic actor-critic line of work developed through deterministic policy gradients and DDPG \citep{silver2014deterministic,lillicrap2015continuous}. It addresses overestimation and function-approximation error through mechanisms such as twin critics, delayed policy updates, and target policy smoothing \citep{fujimoto2018td3}. Together, PPO, SAC, and TD3 provide a baseline set that covers several important design choices in deep reinforcement learning.

GRPO motivates the final project direction, but it is not evaluated as a control method in this interim report. In DeepSeekMath, GRPO is introduced as a PPO-related reinforcement learning method for mathematical reasoning in language models, where the policy update signal is computed relative to groups of sampled outputs \citep{shao2024deepseekmath}. This idea is relevant because it suggests an alternative to conventional value-based advantage estimation. However, transferring this idea to online control is not direct. Control tasks involve sequential environment interaction, temporal credit assignment, continuous or structured action spaces, and evaluation across random seeds and training budgets. It therefore remains open whether GRPO-style group-relative updates can be formulated and evaluated meaningfully for control benchmarks.

The project uses ObjectRL and Gymnasium as part of a reproducible experimental framework. ObjectRL provides the baseline implementations used for the vector-control PPO, SAC, and TD3 experiments, while Gymnasium-style interfaces support standardized interaction with environments \citep{baykal2025objectrl,towers2024gymnasium}. CarRacing is handled through project-side CNN agents because its image observations do not fit directly into the vector-observation workflow used for the ObjectRL-based runs. This tooling choice is not an algorithmic contribution, but it explains how the baseline experiments are organized and validated consistently.

Overall, the related work places the project between established control baselines and recent GRPO-motivated policy optimization from language-model reasoning. The interim report does not answer whether GRPO is better than PPO, SAC, or TD3 for control. Instead, it establishes the necessary comparison basis: a complete and validated PPO/SAC/TD3 baseline matrix on three environments with five seeds. The final project stage will build on this foundation by formulating, implementing, and evaluating a GRPO-control variant under the same experimental structure.
